\id{МРНТИ 31.15.33, 61.53.15}{}

\begin{header}
\swa{О.В.~Рожкова, О.О.~Ковалева, Ж.А.~Сеитова, В.И.~Рожков, В.И Лебедева}{ОСОБЕННОСТИ ПОЛУЧЕНИЯ И СТАБИЛИЗАЦИИ ИНГИБИРУЮЩИХ СВОБОДНО ДИСПЕРСНЫХ ГИДРОГЕЛЕВЫХ ПРОМЫВОЧНЫХ СИСТЕМ КОАГУЛЯЦИОННОГО ТИПА}

\tsp{1,2}О.В.~Рожкова\envelope,
\tsp{3}О.О.~Ковалева,
\tsp{1}Ж.А.~Сеитова,
\tsp{1}В.И.~Рожков,
\tsp{4}В.И Лебедева
\end{header}

\begin{affil}
\tsp{1}Казахский~агротехнический исследовательский университет им.~С.~Сейфуллина, Астана, Казахстан,

\tsp{2}АО «PolyTechPark», Алматы, Казахстан,

\tsp{3}Алтайский геолого-экологический институт, Усть-Каменогорск, Казахстан,

\tsp{4}ООО «Бентонит Хакасии», Черногорск, Россия

\corrauthor{Корреспондент автор: rozhkova.o@stsolutions.kz}
\end{affil}

В работе рассмотрены физико-химические основы формирования, стабилизации
и управляемой пептизации ингибирующих свободнодисперсных гидрогелевых
промывочных систем коагуляционного типа, применяемых в технологии
строительства нефтегазовых скважин. Проведён сравнительный анализ
глинистых и полимерных гидрогелей, выявлены ключевые ограничения их
практического использования, связанные с низкой коагуляционной
устойчивостью бентонитовых суспензий, а также термо- и биодеструкцией
полимерных гелей.

На основе представлений DLVO-теории, коллоидной химии и термодинамики
коагуляционных процессов обоснованы условия получения устойчивых
бентонитовых гидрогелей с пептизированной дисперсной фазой. Показано,
что предварительное разбавление глинистой суспензии до оптимальной
концентрации коллоидных частиц, в сочетании с полимерной лиофилизацией
поверхности дисперсной фазы, позволяет существенно снизить скорость
медленной коагуляции и повысить электростатическую и
структурно-механическую устойчивость системы.

Установлены закономерности влияния анионных и неионных полимеров
полисахаридного ряда на формирование защитных адсорбционно-сольватных
слоёв, величину электростатического барьера коагуляции и глубину
вторичного потенциального минимума. Обоснована целесообразность
комбинированного применения анионных (карбоксиметилцеллюлоза) и неионных
(крахмал, лигносульфонаты) полимеров для обеспечения оптимального
соотношения подвижности и прочности защитных слоёв.

Показано, что введение коагулирующих солей щелочных, щелочноземельных и
амфотерных металлов приводит к развитию обратимой коагуляции во
вторичном потенциальном минимуме за счёт снижения заряда поверхности
глинистых частиц и частичного разрушения полимерных адсорбционных
оболочек. Сформулированы основные правила управляемой пептизации
полимер-глинистых систем, реализуемой механическим или ионным способом,
и определены условия, при которых коагуляция сохраняет обратимый
характер.

Полученные результаты позволяют целенаправленно регулировать
реологические, структурно-механические и ингибирующие свойства
гидрогелевых промывочных систем и могут быть использованы при разработке
высокоэффективных малоглинистых и недиспергирующих буровых растворов для
сложных термобарических условий.

{\bfseries Ключевые слова}: бентонит, гидрогель, полисахарид, пептизация,
коагуляция, полимер-глинис\-тый раствор, лиофильная суспензия.

\begin{header}
КОАГУЛЯЦИЯЛЫҚ ТИПТЕГІ ЕРКІН ДИСПЕРСТІ ГИДРОГЕЛЬДІ ЖҮЙЕЛЕРІН АЛУ ЖӘНЕ ТҰРАҚТАНДЫРУ ЕРЕКШЕЛІКТЕРІ

\tsp{1,2}О.В.~Рожкова\envelope,
\tsp{3}О.О.~Ковалева,
\tsp{1}Ж.А.~Сеитова,
\tsp{1}В.И.~Рожков,
\tsp{4}В.И Лебедева
\end{header}

\begin{affil}
\tsp{1}С.Сейфуллин атындағы Қазақ агротехникалық зерттеу университеті, Астана, Қазақстан,

\tsp{2}«PolyTechPark» АҚ, Алматы, Қазақстан,

\tsp{3}Алтай геологиялық-экологиялық институты, Өскемен, Қазақстан,

\tsp{4}«Бентонит Хакасия» ЖШҚ,Черногорск, Ресей,

e-mail: rozhkova.o@stsolutions.kz
\end{affil}

Жұмыста мұнай-газ ұңғымаларын салу технологиясында қолданылатын
коагуляциялық типтегі ингибирлеуші еркін-дисперсті гидрогельді жуу
жүйелерін қалыптастырудың, тұрақтандырудың және басқарылатын
пептизациялаудың физика-химиялық негіздері қарастырылған. Балшықты және
полимерлі гидрогельдерге салыстырмалы талдау жүргізіліп, бентонитті
суспензиялардың коагуляциялық тұрақтылығының төмендігімен, сондай-ақ
полимерлі гельдердің термо- және биодеструкцияға ұшырауымен байланысты
олардың практикалық қолданылуын шектейтін негізгі кемшіліктер
анықталған.

DLVO-теориясының, коллоидтық химияның және коагуляциялық процестер
термодинамикасының қағидалары негізінде пептизирленген дисперсті фазасы
бар тұрақты бентонитті гидрогельдерді алудың шарттары негізделген.
Балшықты суспензияны коллоидтық бөлшектердің оңтайлы концентрациясына
дейін алдын ала сұйылту және дисперсті фазаның бетін полимерлік
лиофилизациялау баяу коагуляция жылдамдығын едәуір төмендетіп, жүйенің
электростатикалық және құрылымдық-механикалық тұрақтылығын арттыратыны
көрсетілген.

Полисахаридтер қатарының аниондық және иондық емес полимерлерінің
қорғаныш адсорбция\-лық-сольваттық қабаттардың түзілуіне, коагуляцияның
электростатикалық тосқауылының биіктігіне және екіншілік потенциалдық
минимумның тереңдігіне әсер ету заңдылықтары анықталған. Қорғаныш
қабаттардың қозғалғыштығы мен беріктігінің оңтайлы арақатынасын
қамтамасыз ету үшін аниондық (карбоксиметилцеллюлоза) және иондық емес
(крахмал, лигносульфонаттар) полимерлерді біріктіріп қолданудың мақсатқа
сай екендігі негізделген.

Сілтілік, сілтілік-жер және амфотерлік металдардың коагуляциялаушы
тұздарын енгізу глинистi бөлшектер бетінің зарядының төмендеуі мен
полимерлі адсорбциялық қабықшалардың ішінара бұзылуы нәтижесінде
екіншілік потенциалдық минимум аймағында қайтымды коагуляцияның дамуына
әкелетіні көрсетілген. Полимер-балшықты жүйелердің басқарылатын
пептизациясының механикалық және иондық жолмен жүзеге асырылу ережелері
тұжырымдалып, коагуляцияның қайтымды сипаты сақталатын жағдайлар
анықталған.

Алынған нәтижелер гидрогельді жуу жүйелерінің реологиялық,
құрылымдық-механикалық және ингибирлеуші қасиеттерін мақсатты түрде
реттеуге мүмкіндік береді және күрделі термобарикалық жағдайлар үшін
жоғары тиімді азбалшықты және диспергирленбейтін бұрғылау ерітінділерін
әзірлеуде қолданылуы мүмкін.

{\bfseries Түйін сөздер:} бентонит, гидрогель, полисахарид,
пептизация, коагуляция, полимер-балшықері\-тіндісі, гидролиофильді
суспензия.

\begin{header}
FEATURES OF OBTAINING AND STABILIZING INHIBITORY FREE-DISPERSED HYDROGEL COAGULATION-TYPE FLUSHING SYSTEMS

\tsp{1,2}O.V.~Rozhkova\envelope,
\tsp{3}O.O.~Kovaleva,
\tsp{1}Zh.A.~Seitova,
\tsp{1}V.I.~Rozhkov,
\tsp{4}V.I.~Lebedeva
\end{header}

\begin{affil}
\tsp{1}S. Seifullin Kazakh Agrotechnical Research University, Astana, Kazakhstan,

\tsp{2}JSC «PolyTechPark», Almaty, Kazakhstan,

\tsp{3}Altai Geological and Ecological Institute, Ust-Kamenogorsk, Kazakhstan,

\tsp{4}LLC «Bentonite of Khakassia», Chernogorsk, Russia,

e-mail: rozhkova.o@stsolutions.kz
\end{affil}

The paper examines the physicochemical fundamentals of the formation,
stabilization, and controlled peptization of inhibitory free - dispersed
hydrogel flushing systems of the coagulation type used in oil and gas
well construction. A comparative analysis of clay - based and polymer
hydrogels was carried out, and the key limitations of their practical
application were identified, associated with the low coagulation
stability of bentonite suspensions as well as the thermal and
biodegradation of polymer gels.

Based on the concepts of DLVO theory, colloid chemistry, and the
thermodynamics of coagulation processes, the conditions for obtaining
stable bentonite hydrogels with a peptized dispersed phase were
substantiated. It was shown that preliminary dilution of the clay
suspension to an optimal concentration of colloidal particles, combined
with polymer-induced lyophilization of the dispersed phase surface,
signifi\-cantly reduces the rate of slow coagulation and enhances the
electrostatic and structural - mechanical stability of the system.

The regularities of the influence of anionic and nonionic
polysaccharide-based polymers on the formation of protective adsorption
- solvation layers, the height of the electrostatic coagulation barrier,
and the depth of the secondary potential minimum were established. The
expediency of the combined use of anionic (carboxymethyl cellulose) and
nonionic (starch, lignosulfonates) polymers to ensure an optimal balance
between the mobility and strength of protective layers was
substantiated.

It was demonstrated that the introduction of coagulating salts of
alkali, alkaline earth, and amphoteric metals leads to the development
of reversible coagulation in the secondary potential minimum due to a
decrease in the surface charge of clay particles and partial destruction
of polymer adsorption shells. The basic rules for the controlled
peptization of polymer - clay systems, implemented by mechanical or
ionic methods, were formulated, and the conditions under which
coagulation retains a reversible character were determined.

The obtained results make it possible to purposefully regulate the
rheological, structural -mechanical, and inhibitory properties of
hydrogel flushing systems and can be used in the development of highly
efficient low-clay and non - dispersing drilling fluids for complex
thermobaric conditions.

{\bfseries Keywords:} bentonite, hydrogel, polysaccharide, peptization,
coagulation, polymer-clay solution, lyo\-philic suspension.

\begin{multicols}{2}
{\bfseries Введение}. Гидрогелевые промывочные жидкости находят широкое
практическое применение в технологии строительства нефтегазовых скважин.
Данные системы отличаются способом получения и природой
структурообразователя: полимерные гидрогели, глинистые гидрогели,
солевые гидрогели. Полимерные и глинистые гидрогели относятся к
растворам с диспергированной фазой, солевые гидрогели - коллоидные
системы с конденсированной фазой.

В технологии получения структурированных буровых растворов традиционным
является диспергационный способ. Основным структурообразователем
является глина, которая после ввода гидратируется, диспергируется и
создает коагуляционную структуру суспензии. В такой суспензии
удерживаются другие диспергированные материалы - утяжелители, выбуренный
шлам.

Минеральные гидрогели существенно дешевле полимерных аналогов, так как
производятся из крупнотоннажных продуктов переработки и обогащения
горных пород. В практике бурения наиболее распространенным минеральным
гидрогелем является бентонитовый буровой раствор, использующий эффект
набухания и диспергирования в водной среде слоистого алюмосиликата
смектитового типа. Такого рода минеральные гидрогели можно отнести к
системам с пептизированной дисперсной фазой.

Несмотря на все преимущества бентонитовых растворов (псевдопластические
реологические характеристики, оптимальное соотношение цена/качество,
возможность применения малоглинистых ингибирующих и недиспергирующих
систем и др.), тем не менее, они обладают рядом недостатков, существенно
ограничивающих их применение: высокий уровень адгезии глинистой
поверхности к металлу и породам, что приводит к образованию сальников;
низкая коагуляционная устойчивость глинистой суспензии в условиях
притока пластовых вод и под действием высоких температур; большие объемы
наработки при бурении в гигроскопичных глинах; относительно низкое
качество крепления скважин обсадными колоннами в результате
контракционного изменения остатков глинистой фильтрационной корки.

{\bfseries Материалы и методы.} Гидрогели на основе полимеров имеют
вязкопластический характер течения и обратимое восстановление
реологических характеристик при циркуляции раствора в кольцевом
пространстве {[}1-4{]}. В водно-полимерных гелях основными
структурно-механическими единицами являются макромолекулы, участвующие в
построении единой пространственной структурной сетки за счет
донорно-акцепторных или ионных связей (рис.1). Примером ионного
механизма формирования полимерного гидрогеля является координация
функциональных групп полиэлектролитов (КМЦ, ГПАА, гипан) вокруг катионов
поливалентных металлов {[}5,6{]}. Прочность таких гелей возрастает с
увеличением мольной доли катиона-комплексообразователя. В работе {[}7{]}
представлен механизм образования пространственно-сшитых полимерных
структур за счет подандных взаимодействий между функциональными группами
полианионов (КМЦ, ПАЦ) и полиэфирными цепочками неионных ПАВ
(полиоксиэтилен (ПОЭ), неонол, проксанол и др.).
\end{multicols}

\fig[0.5\textwidth]{c2/image15}[Рис.1 - Поперечно-сшитый полимерный гидрогель]

\begin{multicols}{2}
В этом случае полиоксиэтиленовые группы ПАВ образуют с противоионами
(Na\tsp{+}) ассоциированный поликатион, способный к
формированию пространственной структуры за счет интерполимерного
комплексообразования с макромолекулами анионного стабилизатора.

В водных растворах ксантанового биополимера формируются водородные связи
между соседними макромолекулами, отличающимися регулярным чередованием
разветвлений трисахаридов (рис.2).

Это обеспечивает оптимальное соотношение структурных и реологических
свойств биополимерных гелей, которые, однако, весьма неустойчивы к
термобарическим воздействиям {[}8{]}. Повышенной прочностью структуры
обладают полимер-дисперсные гели, образованные неионогенными полимерами
(ОЭЦ, ГЭЦ, ПОЭ) и солями переходных и амфотерных элементов (Fe, Ni, Al)
{[}5, 6{]}.

Накопленные данные о физико-химических и технологических свойствах
полимерных гидрогелей с различной природой пространственного
структурообразования позволяют утверждать, что эти системы сочетают
высокие структурные и вязко-пластичные свойства, а также значительные
ингибирующие и смазочные характеристики, свойственные насыщенным
эмульсионным системам.
\end{multicols}

\fig{c2/image16}[Рис.2 - Гидрогель с разветвленной структурой базового полимера, образующего систему водородных связей]

\begin{multicols}{2}

Тем не менее, полимерные гидрогели имеют несколько существенных
недостатков, связанных с особенностями их коллоидно-химического
строения.

Во-первых, прочность структуры полимерного гидрогеля изменяется симбатно
пластической вязкости (η), поэтому в такой системе практически
невозможно достичь псевдопластической реологии. Иначе говоря, величина
пластической вязкости многократно превосходит первое значение
статического напряжения сдвига (η\textgreater\textgreater{}
СНС\tsb{1}), что приводит к недостаточной выносящей
способности бурового раствора и повышенным гидравлическим сопротивлениям
{[}9{]}.

Во-вторых, время жизни полимерного гидрогеля ограничено устойчивостью
базового полимера к био- и термодеструкции. Кроме того, данные гели
хорошо смачивают поверхность смектитовыхминералов, что приводит к
большой наработке бурового раствора при бурении в активных глинах.

{\bfseries Обсуждение и результаты.} Для преодоления недостатков,
характерных для глинистых и полимерных гидрогелей (высокая
гидрофильность и низкая коагуляционная устойчивость
бентонитовыхгидрогелей, термо- и биодеструкция полимерных гелей) были
разработаны и исследованы свободнодисперсные системы ингибирующего типа,
содержащие в рецептуре катионы калия, кальция или алюминия.

Перед введением солей-ингибиторов необходимо обеспечить максимальную
устойчивость глинистой (бентонитовой) суспензии. Повышение устойчивости
гидрозоля состоит из следующих этапов:

- снижение скорости медленной коагуляции системы;

- максимальная лиофилизация частиц дисперсной фазы.

На первом этапе бентонитовую суспензию, используемую в качестве бурового
раствора для бурения вышележащих интервалов, разбавляют водой до
условной вязкости 30 -- 35 с. При этом, достигаются две цели: снижается
концентрация коллоидных частиц (менее 5 мкм) до 20 - 25
кг/м\tsp{3}, а также снижается концентрация присутствующих в
системе водорастворимых солей (хлориды и гидрокарбонаты натрия).

Необходимость снижения концентрации частиц дисперсной фазы продиктована
тем, что время Z, в течение которого количество частиц в результате
коагуляции уменьшится в 2 раза, обратно пропорционально исходной
концентрации дисперсных частиц (1):

\begin{equation}
Z = \frac{3 \eta}{8 R T \upsilon_{o} \alpha}, \quad (1)
\end{equation}

где \\
$Z$ --- время, в течение которого половина частиц скоагулирует; \\
$k$ --- константа Больцмана; \\
$T$ --- абсолютная температура; \\
$\upsilon_{o}$ --- исходная концентрация коллоидных частиц глины; \\
$\eta$ --- динамическая вязкость дисперсионной среды; \\
$\alpha$ --- доля активных соударений частиц, сопровождающихся актом коагуляции.

Кроме того, в результате разбавления гидрозоля снижается концентрация С
1-1-валентного электролита в дисперсионной среде. При этом, толщина
диффузного слоя, предохраняющего частицы дисперсной фазы от коагуляции,
увеличивается, что приводит к повышению электрокинетического потенциала
(потенциал, обуславливающий взаимное отталкивание дисперсных частиц) и,
как следствие, повышается устойчивость суспензии.

Признаком достаточности разбавления, наряду со снижением условной
вязкости, является уменьшение статического напряжения сдвига до 5/9
Фунт/100Фут\tsp{2}. Это означает, что суспензия находится в
состоянии медленной коагуляции (частицы размещены в потенциальных ямах,
и система характеризуется дальним порядком), однако, Ван-дер-Ваальсовые
силы притяжения между соседними частицами находятся на достаточно
безопасном низком уровне для коагуляционной системы. Тиксотропия золя
(отношение СНС\tsb{10}/СНС\tsb{1}), характеризующая
скорость медленной коагуляции, для разбавленной суспензии не должна
превышать 2.

Важной характеристикой готовности суспензии к введению коагулирующей
соли является пластическая вязкость, которая не должна превышать 10 --
13 мПа·с. Это означает, что глинистые дисперсные частицы имеют
минимальный размер, и в системе отсутствует коагуляция.

Целью полимерных обработок предварительно разбавленной суспензии
являются:

- повышение отрицательного заряда поверхности дисперсных частиц глины;

- увеличение толщины диффузной оболочки, окружающей коллоидные частицы;

- обеспечение сорбционно-сольватного и структурно-механического факторов
устойчивости глинистой суспензии;

- аккумуляция коагулирующих катионов в составе водорастворимых
ион-дипольных комплексов, существующих в дисперсионной среде бурового
раствора.

Введение в суспензию полимерных реагентов, макромолекулы которых имеют
отрицательный заряд (одноименный с зарядом поверхности дисперсных
глинистых частиц), таких как, карбоксиметилированные производные
целлюлозы, а также лигносульфонаты, приводит к появлению вокруг
коллоидных частиц защитных адсорбционных слоев. При этом, эффективный
отрицательный заряд поверхности дисперсных частиц увеличивается, и,
соответственно, повышается энергия отталкивания соседних дисперсных
частиц, изменяющаяся в соответствии с уравнением Дерягина-Ландау (2):

\begin{equation}
U = 64 \left( \frac{c R T}{\kappa} \right) W^{2} e^{-\kappa h} - \frac{A}{h^{2}},
\end{equation}

где \\
$c$ --- концентрация электролита в дисперсионной среде; \\
$\kappa$ --- величина, обратная толщине диффузной оболочки; \\
$W$ --- параметр, пропорциональный отрицательному потенциалу поверхности коллоидной глинистой частицы; \\
$h$ --- расстояние между соседними коллоидными частицами; \\
$A$ --- константа Гамакера, характеризующая Ван-дер-Ваальсово притяжение между соседними коллоидными частицами.

Согласно данному уравнению, на дальних расстояниях между коллоидными
частицами преобладает притяжение (вторичный потенциальный минимум
(потенциальная яма)). На средних расстояниях между коллоидными
глинистыми частицами преобладает отталкивание (электростатический барьер
коагуляции). На ближних расстояниях, при максимальном сближении
коллоидных частиц, преобладает притяжение -- первичная потенциальная
яма.

Возможность коагуляции определяется высотой барьера и глубиной
потенциальных ям (рис.3).
\end{multicols}

\fig[0.5\textwidth]{c2/image17}[\normalfont{\emph{1 - Первичный потенциальный минимум
(яма); 2 - Электростатический барьер отталкивания; 3 - Вторичный
потенциальный минимум (яма)}}\\{\bfseries Рис.3 - Зависимость полной энергии
взаимодействия коллоидных частиц от расстояния (r) между ними}]

\begin{multicols}{2}
Для глинистой суспензии электростатический барьер достаточно велик
(значительно превышает величину теплового движения частиц k∙T), а
вторичный потенциальный минимум также много больше k∙T. Происходит
дальнее взаимодействие двух коллоидных частиц, фиксируемых на расстоянии
порядка 100 нм, отвечающем вторичной потенциальной яме. При этом
наблюдается акт медленной коагуляции, и к данным частицам присоединяются
также на дальних расстояниях другие частицы, что и приводит к появлению
структурированного золя, обладающего статическим напряжением сдвига.

Таким образом, величина СНС косвенно характеризует энергию
Ван-дер-Ваальсового притяжения частиц (глубину вторичной потенциальной
ямы), а тиксотропия суспензии
(СНС\tsb{10}/СНС\tsb{1}) - скорость медленной
коагуляции во вторичной потенциальной яме.

Таким образом, введение отрицательно заряженных полимеров повышает
величину параметра W пропорциональную заряду дисперсной частицы, что,
соответственно, приводит к повышению энергии отталкивания между
частицами, уменьшению глубины вторичной потенциальной ямы, что
проявляется в снижении величины СНС суспензии (рис.4). Вряду 1÷3
увеличивается концентрация анионного полимера и, соответственно,
увеличивается высота электростатического барьера отталкивания и
снижается глубина второго потенциального мини-мума.

Отсюда следует, что полимерные обработки должны приводить к снижению СНС
глинистой суспензии.

При введении в глинистую суспензию, как заряженных, так и неионных
полимеров (крахмал, оксиэтилцеллюлоза и пр.) происходит формирование
полимерных адсорбционных слоев на поверхности частиц дисперсной фазы, за
счет которых толщина диффузной оболочки коллоидной частицы
увеличивается. При этом происходит рост энергии отталкивания частиц в
области второго потенциального минимума и повышается коагуляционная
устойчивость системы (рис.5). Вряду 1÷4 увеличивается концентрация
неионного полимера.
\end{multicols}

\begin{figs}
    \fig[0.45\textwidth]{c2/image18}[Рис.4 - Влияние концентрации
анионных полимеров на полную энергию взаимодействия коллоидных частиц]
\fig[0.45\textwidth]{c2/image19}[Рис.5 - Влияние концентрации
неионогенных полимеров на полную энергию взаимодействия коллоидных
частиц]
\end{figs}

\begin{multicols}{2}
Обеспечение сорбционно-сольватного и структурно-механического факторов
устойчивости глинистой суспензии осуществляется, во-первых, за счет
введения полимеров в глинистую суспензию, которое существенно повышает
динамическую вязкость дисперсионной среды, что приводит к замедлению
коагуляционных процессов (полупериод коагуляции возрастает). Во-вторых,
адсорбция полимеров на поверхности коллоидных глинистых частиц приводит
к образованию весьма развитых сольватных слоев. Большие размеры
макромолекул полимеров, несущих собственные сорбционно-сольватные слои
значительной протяженности и прочности, создают на поверхности
коллоидных частиц адсорбционно-сольватные слои значительной плотности,
которые перекрывают не только первичную, но и вторичную потенциальную
яму. Устойчивость таких лиофилизированных суспензий весьма высока. Таков
механизм формирования структурно-механического барьера коагуляции.

Прочность полимерных адсорбционных слоев увеличивается во времени,
достигая предельного значения лишь через несколько часов, что связано с
замедленной диффузией макромолекул и медленной их ориентацией на
поверхности коллоидных частиц. Поэтому, после введения в глинистую
суспензию полимерной обработки, следует в течение достаточно длительного
времени (минимум 1 -- 1,5 часа) перемешивать суспензию в блоке
приготовления раствора. Коагулирующие соли вводят только через 1,5 -- 2
ч. после того, как была произведена полимерная стабилизация глинистой
суспензии.

Следует отметить, что стабилизацию глинистых суспензий наиболее часто
производят анионными полимерами. При этом, влияние электрических зарядов
функциональных групп полимера на прочность адсорбционных слоев
оказывается весьма сложным. Так, повышение заряда (в результате
диссоциации функциональных групп макромолекул), усиливающее гидратацию,
создает отталкивание между макромолекулами, уменьшая боковую когезию, а,
следовательно, и прочность адсорбционного слоя. В ряду производных
полисахаридов прочность адсорбционных слоев повышается в ряду: ПАЦ → КМЦ
→ Крахмал → Гидроксиэтилцеллюлоза.

Поэтому, стабилизацию глинистой суспензии перед введением коагулирующих
солей следует осуществлять за счет обработок КМЦ и солестойким
крахмалом. Следует также поддерживать не слишком высокий уровень рН
системы раствора (не более 9), при котором степень диссоциации
карбоксильных групп полимера будет существенно ниже 1.

Недостатком неионных полимеров (крахмала) является слишком высокая
прочность насыщенного адсорбционного слоя, что приводит, в конце концов,
к резкому снижению стабилизирующего действия. Это происходит вследствие
того, что случайные разрывы сплошности адсорбционного слоя уже не могут
«самозалечиваться» из-за «хрупкости» сверхпрочного адсорбционного слоя,
который становится малоподвижным. Поэтому, при стабилизации глинистой
суспензии оптимальной стратегией является сочетание анионных (КМЦ) и
неионных полимеров. Анионные полимеры обеспечивают необходимую
подвижность и мобильность защитных адсорбционных слоев, а неионные
полимеры создают необходимую прочность адсорбционных слоев {[}9-11{]}.

Наряду с высокой прочностью и вязкостью, полимерные адсорбционные слои
обеспечивают устойчивость глинистых золей также за счет других факторов.
Так, в зазоре между двумя частицами, окруженными сорбционно-сольватными
слоями, происходит увеличение концентрации, что приводит к появлению
осмотического расклинивающего давления, препятствующего коагуляции.
Кроме того, коагуляция частиц, покрытых полимерными адсорбционными
слоями запрещена вторым законом термодинамики, согласно которому
энтропия системы стремится к максимуму. При сближении частиц, окруженных
полимерными слоями, происходит резкое понижение энтропии вследствие
ограничения свободы движения сегментов полимерных цепей при перекрытии
адсорбционных слоев.

Применяемые для стабилизации глинистых суспензий полимерные реагенты
должны обладать способностью к образованию водорастворимых ион-дипольных
комплексов с коагулирующими катионами. Это обеспечивает снижение
адсорбции катионов на поверхности дисперсных частиц, а также существенно
повышает ингибирующий эффект полимер-глинистых растворов. К такого рода
полимерам относятся производные полисахаридов -- КМЦ, ПАЦ, крахмалы.
Кроме того, коагуляционную устойчивость глинистых растворов существенно
повышают полимеры, обладающие эфирным кислородом - полиоксиэтилен (ПОЭ),
полиалкиленгликоли (сополимеры окисей этилена и пропилена),
полипропиленгликоль (ППГ) и пр.

{\bfseries Выводы.} Таким образом, при стабилизации глинистой суспензии
перед введением коагулирующих солей необходимо следовать нескольким
правилам:

1. Разбавление базового глинистого раствора, переведенного с предыдущего
интервала, до концентрации коллоидных частиц (MBT) не более 20 -- 25
кг/м\tsp{3}. Так как распределение коллоидных частиц по
размерам для разных суспензий может сильно различаться, то вторичными
признаками готовности суспензии для введения коагулирующей соли являются
следующие значения технологических свойств: СНС=5/9
Фунт/100Фут\tsp{2}, η\tsb{пл.}<13 мПа∙с,
УВ<35 с, СНС\tsb{10}/СНС\tsb{1} <{}
2.2. Введение полимерных реагентов для лиофилизации поверхности дисперсной
фазы. Обязательным является введение полимеров в глинистую суспензию в
виде водных растворов. Необходимо использовать два типа стабилизирующих
полимеров -- спирально ориентированные анионные (КМЦ) полимеры и
гибкоцепные линейные либо неионные полимеры - лигносульфонаты, крахмалы.
Не рекомендуется использовать сильнозарядные полимеры (ПАЦ), а также
линейные полимеры с повышенной гибкостью углеводородной цепи (акрилаты,
полиакриламиды). Первые из них (ПАЦ) из-за высокого заряда макромолекул
образуют слишком разряженные и непрочные адсорбционные слои, а вторые
(ГПАА) обладают повышенной гидрофобностью и гибкостью углеводородной
матрицы (отсутствует эфирный кислород), что приводит к низкой
устойчивости к действию коагулирующих катионов. Расход полимеров на
предварительную стабилизацию глинистой суспензии составляет (в пересчете
на 1 кг коллоидной глинистой фазы): 0,3 -- 0,4 кгКМЦ; 0,7 - 0,8
кггибкоцепных лигносульфоната или крахмала.

3. После введения полимеров в глинистую суспензию следует произвести
тщательное перемешивание раствора в течение не менее 1,5 часов для
обеспечения защитных адсорбционных слоев требуемой прочности.

4. Уровень рН полимер-глинистого раствора не следует повышать выше 9 --
9,5, чтобы не спровоцировать слишком большое повышение заряда полимерных
макромолекул КМЦ, снижающее прочность защитных адсорбционных слоев.

После введения коагулирующей соли (алюмокалиевые квасцы, гипс, хлористый
кальций и пр.) наблюдается увеличение вязкости полимер-глинистого
раствора, а также повышение его структурных свойств (рост СНС). Причиной
этого является повышение межмолекулярных Ван-дер-Ваальсовых сил,
усиливающее притяжение между соседними частицами дисперсной глинистой
фазы. То есть, происходит развитие коагуляции во вторичной потенциальной
яме вследствие увеличения ее глубины.

Усиление притяжения между коллоидными частицами (углубление вторичной
потенциальной ямы) происходит из-за того, что коагулирующие катионы
(Са\tsp{2+}) образуют нерастворимые соединения с
карбоксилатными группами макромолекул полимеров, образующих защитный
адсорбционный слой. То есть, в адсорбционно-сольватных слоях возникают
«бреши». Кроме того, коагулирующие катионы (Са\tsp{2+},
К\tsp{+}) обладают способностью к нейтрализации поверхности
глинистых коллоидных частиц в результате сорбции на их поверхности. Это
приводит к снижению отрицательного заряда (и потенциала) поверхности
дисперсных частиц, что снижает электростатическое отталкивание
(электростатический «барьер») и увеличивает глубину вторичной
потенциальной ямы.

При введении коагулирующих солей заряд коллоидных частиц уменьшается,
прочность защитных полимерных адсорбционных слоев снижается, частицы
подходят ближе друг к другу и, соответственно, развивается коагуляция
системы во вторичном потенциальном минимуме. Из-за этого повышаются
реологические и структурные характеристики полимер-глинистого раствора.

Если глинистая суспензия перед введением коагулирующих солей была как
следует лиофилизирована полимерными обработками, то коагуляция во
вторичном потенциальном минимуме носит обратимый характер, то есть,
коагулированная система обладает способностью к пептизации.

Пептизация тем более вероятна, чем более лиофилизирована (обработана
полимерами) глинистая суспензия и чем меньше времени прошло с момента
коагуляции. С течением времени при ближнем взаимодействии происходит
постепенное срастание частиц с уменьшением дисперсности и поверхностной
энергии. В этом случае коагуляция приобретает необратимый характер, и
пептизация исключается.

Различают два способа пептизации полимер-глинистой суспензии:

1. Механический способ. Данный метод заключается в том, что
полимер-глинистая суспензия, загустевшая после введения коагулирующих
катионов, должна быть, как можно быстрее деформирована путем
интенсивного перемешивания с применением механических и гидравлических
средств. Время перемешивания зависит от его интенсивности и составляет
от 10 до 30 мин. Перемешивание может заменить циркуляция
скоагулировавшего полимер-глинистого раствора через скважину в течение 1
- 2 циклов. При этом суспензия претерпевает весьма высокие деформации
при прохождении через насадки гидромониторного долота.

Принцип механической пептизации заключается в том, что коллоидные
дисперсные частицы, за счет приложенных сдвиговых усилий, смещаются друг
относительно друга. Причем расстояние между соседними частицами
увеличивается. В этот момент «бреши» в адсорбционных полимерных слоях
затягиваются (происходит «самозалечивание» адсорбционных слоев) в
результате диффузии макромолекул полимеров в направлении обедненных
участков адсорбционного слоя. В результате, после прекращения действия
сдвиговых деформаций, коллоидные частицы уже не могут подойти друг к
другу на близкое расстояние - проявляется отталкивающее действие
защитных адсорбционно-сольватных оболочек. При этом, глубина вторичной
потенциальной ямы уменьшается на несколько k∙T, что приводит к снижению
СНС и разжижению полимер-глинистого раствора.

2. Если лиофилизация суспензии произведена недостаточно тщательно
(система обеднена полимерами) или концентрация коллоидной глинистой фазы
повышена (МВТ \textgreater35 -- 40 кг/м\tsp{3}), то для
осуществления пептизации скоагулировавшегополимер-глинистого раствора
следует увеличить заряд и потенциал коллоидной глинистой частицы путем
добавления электролита, содержащего потенциалообразующие
(неиндифферентные) ионы. Для глинистой суспензии такими ионами являются
анионы (отрицательные ионы), адсорбирующиеся на ребрах и сколах
кристаллической решетки. Источниками таких анионов являются, например,
НТФ, ТПФН, хроматы и бихроматы. Концентрация данных анионов в
полимер-глинистом растворе для осуществления пептизации не должна
превышать 0,1 -- 0,3 кг/м\tsp{3}.

Таким образом, механический и ионный механизмы пептизации по своей
природе являются единым процессом, основанным на увеличении энергии
отталкивания между коллоидными частицами за счёт повышения заряда и
потенциала их поверхности, а также увеличения толщины диффузной
оболочки, что приводит к дезагрегации глинистых частиц. Эффективная
стабилизация глинистой суспензии перед введением коагулирующих солей
достигается посредством её предварительного разбавления до оптимальной
концентрации дисперсной фазы и последующей обработки сочетанием анионных
полимеров (КМЦ) и неионных полимеров (крахмал,
лигносульфонаты). Такое комплексное введение полимеров
обеспечивает формирование прочных, но подвижных адсорбционно-сольватных
слоёв, увеличивает электростатический и структурно-механический барьер,
снижает склонность системы к коагуляции и формирует контролируемую,
преимущественно обратимую структуру полимер-глинистого раствора.
\end{multicols}

\begin{center}
{\bfseries Литература}
\end{center}

\begin{refs}
1. Рожкова О.В., Ковалева О.О., Лебедева В.И., Ветюгов Д.А., Рожков
В.И.. Исследование процесса окомкования высококачественного
железорудного концентрата с применением новых связующих добавок //
Вестник КазУТБ. -2025. -Т.2(27). -С.231-245. DOI
10.58805/kazutb.v.2.27-954.

2. Ibraimova D.M-K., Rozhkova O.V., Musabekov K.B., Tazhibayeva S.M.,
Rozhkov V.I., Yermekov M.T. Development of methods to obtain composite
materials from organoclays // Eurasian Journal of Chemistry. -2023.
-Vol.28. -No.4 (112). -P.101-111. DOI 10.31489/2959-0663/4-23-14.

3. Рожкова О.В., Муздыбаева Ш.А., Мусабеков К.Б., Ибраимова. Д.М.-К.,
Рожков В.И., Ермеков М.Т.. Разработка методов получения носителей
лекарственных средств на основе органомодифицированных глин//~Известия
НАН РК. Серия химических наук. -2023. -T.3: 456. -C.138-156. DOI
10.32014/2023.2518-1491.183.

4. Musabekov K.B., Rozhkova O.V., Artykova (Ibraimova) D.M-K., Yermekov
M.T., Muzdybaeva Sh.A. Application of bentonite clay as a protective
barrier in the disposal of radioactive waste from the nuclear industry
of Kazakhstan // News of the National Academy of Sciences of the
Republic of Kazakhstan. Series: Chemistry and Technology. -2023. -Vol.
1(454). -P.66-77. DOI 10.32014/2023.2518-1491.148.

5. Кошелев В.Н., Вахрушев Л.П., Беленко Е.В., Лушпеева О.А.
Полимер-дисперсные синергетические явления и новые системы буровых
растворов // Нефтяное хозяйство. -2001. - № 4. -С.22-23.

6. Беленко Е.В. Общие технологические принципы промывки скважин//
Вестник Ассоциации буровых подрядчиков. -2010. - № 3. - С.2-7.

7. Пеньков А. И., Кошелев В. Н., Вахрушев Л. П., Беленко Е. В., Лушпеева
О. А. Применение термодинамики фазовых переходов для расчётов
технологических параметров буровых растворов // Нефтяное хозяйство. -
2001. - № 9. - С.44-48.

8. Пеньков А.И., Шишов В. А., Федосов Р.И. Требования к растворам
безглинистым и с низким содержанием твердой фазы // Нефтяное хозяйство.
-1981. - № 5. - С.23.

9. Федосов Р. И., Пеньков А. И. Разработка экологически совместимых с
окружающей средой буровых растворов для бурения и закачивания скважин на
арктическом шельфе/ Вопросы промывки скважин с горизонтальными участками
ствола. Краснодар. -1998. -C.49-53.

10. Беленко Е.В., Мазыкин С.В., Мнацаканов В.А. Принципы нормирования
реологических и структурно-механических свойств утяжеленных
полимерглинистых буровых растворов // Строительство нефтяных и газовых
скважин на суше и море. - 2011. - № 7. - С.40-43.
\end{refs}

\begin{center}
{\bfseries References}
\end{center}

\begin{refs}
1. Rozhkova O.V., Kovaleva O.O., Lebedeva V.I., Vetjugov D.A., Rozhkov
V.I.. Issledovanie processa okomkovanija vysokokachestvennogo
zhelezorudnogo koncentrata s primeneniem novyh svjazujushhih dobavok //
Vestnik KazUTB. -2025. -T.2(27). -S.231-245. DOI
10.58805/kazutb.v.2.27-954.{[}in Russian{]}

2. Ibraimova D.M-K., Rozhkova O.V., Musabekov K.B., Tazhibayeva S.M.,
Rozhkov V.I., Yermekov M.T. Development of methods to obtain composite
materials from organoclays // Eurasian Journal of Chemistry. -2023.
-Vol.28. -No.4 (112). -P.101-111. DOI 10.31489/2959-0663/4-23-14.

3. Rozhkova O.V., Muzdybaeva Sh.A., Musabekov K.B., Ibraimova. D.M.-K.,
Rozhkov V.I., Ermekov M.T.. Razrabotka metodov poluchenija nositelej
lekarstvennyh sredstv na osnove organomodificirovannyh glin// Izvestija
NAN RK. Serija himicheskih nauk. -2023. -T.3(456). -S.138-156. DOI
10.32014/2023.2518-1491.183. {[}in Russian{]}

4. Musabekov K.B., Rozhkova O.V., Artykova (Ibraimova) D.M-K., Yermekov
M.T., Muzdybaeva Sh.A. Application of bentonite clay as a protective
barrier in the disposal of radioactive waste from the nuclear industry
of Kazakhstan // News of the National Academy of Sciences of the
Republic of Kazakhstan. Series: Chemistry and Technology. -2023. -Vol.
1(454). -P.66-77. DOI 10.32014/2023.2518-1491.148.

5. Koshelev V.N., Vahrushev L.P., Belenko E.V., Lushpeeva O.A.
Polimer-dispersnye sinergeticheskie javlenija i novye sistemy burovyh
rastvorov // Neftjanoe hozjajstvo. -2001. - № 4. -S.22-23. {[}in
Russian{]}

6. Belenko E.V. Obshhie tehnologicheskie principy promyvki skvazhin//
Vestnik Associacii burovyh podrjadchikov. -2 010. - № 3. - S.2-7. {[}in
Russian{]}

7. Pen' kov A. I., Koshelev V. N., Vahrushev L. P.,
Belenko E. V., Lushpeeva O. A. Primenenie termodinamiki fazovyh
perehodov dlja raschjotov tehnologicheskih parametrov burovyh rastvorov
// Neftjanoe hozjajstvo. - 2001. - № 9. - S.48-51. {[}in Russian{]}

8. Pen' kov A.I., Shishov V. A., Fedosov R.I. Trebovanija
k rastvoram bezglinistym i s nizkim soderzhaniem tverdoj fazy //
Neftjanoe hozjajstvo. -1981. - №5. - S.23. {[}in Russian{]}

9. Fedosov R. I., Pen' kov A. I. Razrabotka
jekologicheski sovmestimyh s okruzhajushhej sredoj burovyh rastvorov
dlja burenija i zakachivanija skvazhin na arkticheskom
shel' fe/ Voprosy promyvki skvazhin s
gorizontal' nymi uchastkami stvola. Krasnodar. -1998. -C.
49-53. {[}in Russian{]}

10. Belenko E.V., Mazykin S.V., Mnacakanov V.A. Principy normirovanija
reologicheskih i strukturno-mehanicheskih svojstv utjazhelennyh
polimerglinistyh burovyh rastvorov // Stroitel' stvo
neftjanyh i gazovyh skvazhin na sushe i more. - 2011. - №7. - S.40-43.
{[}in Russian{]}
\end{refs}

\begin{info}
\hspace{1em}\emph{{\bfseries Сведение об авторах}}

Рожкова О.В. - доктор химических наук, профессор, академик
Национальной академии естественных наук Республики Казахстан, НАО
«Казахский агротехнический исследовательский университет им. С.
Сейфуллина», Астана, Казахстан, AO «Science and Technology Solutions»,
Астана, Казахстан, e-mail: rozhkova.o@stsolutions.kz;

Ковалёва О.О. - магистр, ТОО "Алтайский геолого-экологический
институт", Усть-Каменогорск, Казахстан, е-mail:
kovaleva@agikaz.pro\ul{;}

Сеитова Ж.А. -- кандидат технических наук, и.о. ассоцированного
профессора, НАО «Казахский агротехнический исследовательский
университет им. С. Сейфуллина», Астана, Казахстан,
zhadraseitova3110@gmail.com;

Рожков В.И. - кандидат технических наук, НАО «Казахский
агротехнический исследовательский университет им. С. Сейфуллина»,
Астана, Казахстан, e-mail:
\href{mailto:vitalrza1983@gmail.com}{vitalrza1983@gmail.com};

Лебедева В.И. - заместитель генерального директора по стратегическому
развитию, OOO «Бентонит Хакасии», Черногорск, Республика Хакасия,
Россия, e-mail: lebedeva@bentonit.ru.

\hspace{1em}\emph{{\bfseries Information about the authors}}

Rozhkova O.V. - Doctor of Chemical Sciences, Professor, Academician of
the National Academy of Natural Sciences of the Republic of
Kazakhstan, NAO "Saken Seifullin Kazakh AgroTechnical Research
University", Astana, Kazakhstan, JSC "Science and Technology
Solutions", Astana, Kazakhstan, e-mail: rozhkova.o@stsolutions.kz;

Kovaleva O.O. - Master of Nuclear Energy and Thermophysics, LLP "Altai
Geological-Ecological Institute", Ust-Kamenogorsk, Kazakhstan, e-mail:
kovaleva@agikaz.pro;

Seitova Zh.A. - Candidate of Technical Sciences, Acting Associate
Professor, NAO "Saken Seifullin Kazakh AgroTechnical Research
University", Astana, Kazakhstan, e-mail: zhadraseitova3110@gmail.com;

Rozhkov V.I. - Candidate of Technical Sciences, NAO "Saken Seifullin
Kazakh AgroTechnical Research University", Astana, Kazakhstan, e-mail:
\href{mailto:vitalrza1983@gmail.com}{vitalrza1+983@gmail.com};

Lebedeva V.I. - Deputy General Director for Strategic Development, LLC
"Bentonite Khakassia", Chernogorsk, Republic of Khakassia, Russia,
e-mail: lebedeva@bentonit.ru.
\end{info}
